% -------------
% Macros / Snippets usados comunmente
%
% $ := Uso en mathmode
% #N := Argumento N.
% #N! := Argumento opcional N.
% -------------

% Macros de remplazo (comando nuevo) (comando a reemplazar)
\newcommand{\itz}{\textcolor{Variante7}{$\triangleright$}}
\newcommand{\iitz}{\textcolor{Variante8}{$\diamond$}}
\newcommand{\iiitz}{\textcolor{Variante9}{$\circ$}}
\newcommand{\ienum}{\textcolor{Variante7}{\arabic*.}}
\newcommand{\iienum}{\textcolor{Variante8}{\alph*.}}
\newcommand{\iiienum}{\textcolor{Variante9}{\roman*.}}
\newcommand{\spc}{\textvisiblespace}
\newcommand{\qed}{\hfill $\blacksquare$}
\newcommand{\ra}{\rightarrow} % $
\newcommand{\la}{\leftarrow} % $
\newcommand{\Ra}{\Rightarrow} % $
\newcommand{\sii}{\Longleftrightarrow} % $
\newcommand{\R}{\mathbb{R}} % $
\newcommand{\N}{\mathbb{N}} % $
\newcommand{\Z}{\mathbb{Z}} % $
\newcommand{\C}{\mathbb{C}} % $

% Macros de sustitución (comando nuevo) [argumentos] (comando a reemplazar, #N es el argumento N)
\newcommand{\Eq}[1]{Ec.~(\ref{#1})}
\newcommand{\Fig}[1]{Fig.~\ref{#1}}
\newcommand{\Sec}[1]{Sec.~\ref{#1}}
\newcommand{\IT}[1]{\textit{#1}}
\newcommand{\TT}[1]{\texttt{#1}}
\newcommand{\SC}[1]{\textsc{#1}}
\newcommand{\BF}[1]{\textbf{#1}}
\newcommand{\MB}[1]{\mathbf{#1}} % $

% Redefinición al español de nombres o customización (sin estilado)
\renewcommand{\theequation}{\thesection.\arabic{equation}}
\renewcommand{\arraystretch}{1.5}


% Figuras estilizadas con imagen
%   Figuras con entradas para imagen, centradas y con aparición exactamente donde se coloca el comando.
% -------------------
% #1! : Factor de multiplicación para la longitud (ej. 0.5, 0.75, 0.33), "1" por defecto
% #2  : PATH a la imagen, se asume que está en src/assets
% #3  : Pie de página
% #4  : Etiqueta a tag para referenciar
\NewDocumentCommand{\fimage}{ O{} m m m }{
    \begin{figure}[H]
        \Centering
        \includegraphics[width=#1\linewidth]{src/assets/#2}
        \caption{#3}
        \label{#4}
    \end{figure}
}

% Figuras estilizadas multiuso (y con opción de uso para subfigures)
%   Figuras flexibles con estilado centrado.
% -------------------
% #1! : Ubicación, "H" por defecto
% #2  : Contenido de la figura
% #3  : Pie de página
% #4  : Etiqueta a tag para referenciar
\NewDocumentCommand{\figura}{ O{H} m m m }{
    \begin{figure}[#1]
        \Centering
        #2
        \caption{#3}
        \label{#4}
    \end{figure}
}

% Tablas estilizadas
%   Crea un ambiente que crea una tabla con NiceMatrix, con opciones de color, para primer renglón, y para intercalar.
% -------------------
% #1  : Si la función se llama como \tabla* entonces, se solicita que el texto sea de color claro.
% #2  : Bandera !. Si la lleva se coloca un \hline al inicio
% #3  : Color de fondo del primer renglón
% #4  : Color desvanecido de los renglones impares (sin contar el primero)
% #5  : Estilo de columnas
\NewDocumentEnvironment{tabla}{ t* t! O{Variante4} O{Variante4} m }
{
    \begin{Center}
    \begin{NiceTabular}{#5}
    \CodeBefore
        \rowlistcolors{2}{Background, #3!20!Background}
    \Body
        \IfBooleanTF{#2}{\hline}{}
        \rowcolor{#4}
        \IfBooleanTF{#1}{\RowStyle{\color{Foreground}}}{\RowStyle{\color{Background}}}
}
{
    \end{NiceTabular}
    \end{Center}
}